\section{Estimación de casos de uso}
\label{anexo:estimacion-cu}

El tamaño funcional de F.A.M.A. se estimó mediante Puntos de Casos de
Uso (UCP). El cálculo se reconstruyó a partir de los 14 casos de uso del
capítulo~5. Los factores de ajuste se valoraron con base en el diseño, las
iteraciones y el contexto académico del proyecto. Los puntajes
representan juicios de estimación que pueden revisarse si se precisan
los flujos o los antecedentes del equipo.

\subsection{Clasificación de actores}

El peso de cada actor depende de su forma de interacción con F.A.M.A.
Un sistema con una interfaz de programación definida tiene peso 1.
Una interacción por protocolo sin una API definida o por interfaz de
texto tiene peso 2. Un usuario que opera mediante una interfaz gráfica
tiene peso 3. Para este cálculo se incluye Google Cloud Storage (GCS)
como actor secundario por su participación en el intercambio de archivos.
El orquestador, PostgreSQL y el motor de entrenamiento se consideran
componentes internos de la solución.

\begin{table}[!htbp]
    \centering
    \caption{Pesos de los actores de F.A.M.A.}
    \label{tab:ucp-fama-actores}
    \small
    \renewcommand{\arraystretch}{1.12}
    \begin{tabular}{@{}p{0.34\linewidth}p{0.33\linewidth}lc@{}}
        \hline
        \textbf{Actor} & \textbf{Interacción} & \textbf{Tipo} & \textbf{Peso} \\
        \hline
        Investigador & Interfaz web & Complejo & 3 \\
        Administrador de sistemas & Interfaz web & Complejo & 3 \\
        Google Cloud Storage & API de almacenamiento & Simple & 1 \\
        \hline
        \multicolumn{3}{@{}l}{\textbf{Peso de actores (UAW)}} & \textbf{7} \\
        \hline
    \end{tabular}
\end{table}

\subsection{Clasificación de casos de uso}

Se cuenta como transacción cada intercambio funcional diferenciado con
un actor, incluidos los flujos alternativos que requieren una respuesta
distinta. Una solicitud y su respuesta inmediata se cuentan juntas. Las
operaciones internas, como las consultas a PostgreSQL y las épocas del
entrenamiento, se agrupan en la interacción que las origina. Cuando la
especificación del capítulo~5 es resumida, las transacciones se deducen
de las acciones descritas. Cada identificador se pondera una sola vez,
aunque intervengan varios actores.

Un caso simple tiene hasta 3 transacciones y recibe 5 puntos. Un caso
medio tiene entre 4 y 7 transacciones y recibe 10 puntos. Un caso complejo
tiene más de 7 transacciones y recibe 15 puntos según la escala de
Karner (1993).

\begin{table}[!htbp]
    \centering
    \caption{Transacciones y pesos por caso de uso.}
    \label{tab:ucp-fama-casos}
    \small
    \setlength{\tabcolsep}{4pt}
    \renewcommand{\arraystretch}{1.10}
    \begin{tabular}{@{}p{0.18\linewidth}p{0.55\linewidth}clc@{}}
        \hline
        \textbf{Caso} & \textbf{Interacciones consideradas} &
        \textbf{N.º} & \textbf{Tipo} & \textbf{Peso} \\
        \hline
        \texttt{CU\_SES\_01} & Iniciar sesión. Informar credenciales inválidas. & 2 & Simple & 5 \\
        \texttt{CU\_SES\_02} & Consultar el perfil, editar los datos personales, cambiar la contraseña y revisar la actividad. & 4 & Medio & 10 \\
        \texttt{CU\_INV\_01} & Solicitar la sincronización, consultar GCS, descargar audios y consultar el resultado. Informar un fallo de conexión o credenciales. & 5 & Medio & 10 \\
        \texttt{CU\_INV\_02} & Consultar el catálogo y seleccionar un conjunto de datos. & 2 & Simple & 5 \\
        \texttt{CU\_INV\_03} & Configurar e iniciar el entrenamiento, respaldar los pesos en GCS y consultar las métricas. Informar falta de recursos. & 4 & Medio & 10 \\
        \texttt{CU\_INV\_04} & Consultar el historial y revisar las métricas de un modelo. & 2 & Simple & 5 \\
        \texttt{CU\_INV\_05} & Seleccionar el modelo y recuperar sus pesos desde GCS. & 2 & Simple & 5 \\
        \texttt{CU\_INV\_06} & Enviar un audio y obtener su predicción. Informar el rechazo de un formato inválido. & 2 & Simple & 5 \\
        \texttt{CU\_INV\_07} & Abrir el formulario, enviar la corrección y transferir el aporte a GCS. Cancelar la corrección. & 4 & Medio & 10 \\
        \texttt{CU\_ADM\_01} & Listar, crear, editar y deshabilitar cuentas. Informar un correo duplicado. & 5 & Medio & 10 \\
        \texttt{CU\_ADM\_02} & Configurar el nodo y revisar sus recursos de hardware. & 2 & Simple & 5 \\
        \texttt{CU\_ADM\_03} & Actualizar la configuración de credenciales, verificar la conexión y consultar el uso de almacenamiento. & 3 & Simple & 5 \\
        \texttt{CU\_ADM\_04} & Consultar el historial y el detalle de un registro. Informar que un registro está indisponible. & 3 & Simple & 5 \\
        \texttt{CU\_ADM\_05} & Solicitar la detención de un proceso en ejecución. & 1 & Simple & 5 \\
        \hline
        \multicolumn{4}{@{}l}{\textbf{Peso de casos de uso (UUCW)}} & \textbf{95} \\
        \hline
    \end{tabular}
\end{table}

El catálogo contiene nueve casos simples y cinco medios. La complejidad
del procesamiento acústico y del entrenamiento se refleja también en
los factores técnicos, pues una interacción con el usuario puede
desencadenar varias operaciones internas.

\subsection{Factores técnicos y del entorno}

Cada factor recibe una valoración de 0 a 5. Los valores de 0 a 2
representan incidencia baja, de 3 a 4 incidencia intermedia y 5
incidencia esencial. El puntaje de cada fila resulta de multiplicar
el peso de la pauta por la valoración propuesta para el proyecto.

\begin{table}[!htbp]
    \centering
    \caption{Valoración de los factores técnicos.}
    \label{tab:ucp-fama-tecnicos}
    \small
    \setlength{\tabcolsep}{4pt}
    \renewcommand{\arraystretch}{1.10}
    \begin{tabular}{@{}lp{0.26\linewidth}p{0.40\linewidth}ccc@{}}
        \hline
        \textbf{ID} & \textbf{Factor} & \textbf{Justificación} &
        \textbf{Peso} & \textbf{Valor} & \textbf{Puntos} \\
        \hline
        T1 & Sistema distribuido & Coordinación entre el nodo local y el almacenamiento remoto. & 2 & 5 & 10 \\
        T2 & Rendimiento & Predicción interactiva y entrenamiento limitado por memoria y cómputo. & 1 & 5 & 5 \\
        T3 & Eficiencia del usuario & Operación del flujo acústico desde el panel web. & 1 & 4 & 4 \\
        T4 & Procesamiento interno & Extracción de características, entrenamiento e inferencia. & 1 & 5 & 5 \\
        T5 & Reutilización & Componentes compartidos entre dominios, con configuraciones acústicas específicas. & 1 & 4 & 4 \\
        T6 & Facilidad de instalación & Preparación del nodo, servicios y conexión con GCS. & $0{,}5$ & 4 & 2 \\
        T7 & Facilidad de operación & Paneles de operación y mensajes de estado. & $0{,}5$ & 4 & 2 \\
        T8 & Portabilidad & Contenedores y variantes de ejecución en CPU y GPU. & 2 & 3 & 6 \\
        T9 & Facilidad de cambio & Separación entre interfaz, API y predictores. & 1 & 4 & 4 \\
        T10 & Concurrencia & Trabajos en segundo plano y control de acceso a los recursos del nodo. & 1 & 4 & 4 \\
        T11 & Seguridad & Sesiones, roles, permisos y gestión de credenciales. & 1 & 5 & 5 \\
        T12 & Acceso directo de terceros & El alcance estimado contempla la interfaz web y la integración con GCS. & 1 & 0 & 0 \\
        T13 & Capacitación de usuarios & La puesta en marcha contempla instrucción para los roles de operación. & 1 & 3 & 3 \\
        \hline
        \multicolumn{5}{@{}l}{\textbf{Total ponderado técnico (TFactor)}} & \textbf{54} \\
        \hline
    \end{tabular}
\end{table}

La Tabla~\ref{tab:ucp-fama-entorno} utiliza el valor neutro 3 en los
factores del equipo con información insuficiente para asignar un nivel
más preciso. Las valoraciones se aplican al equipo en conjunto y
pueden actualizarse si se dispone de registros de dedicación o
experiencia previa.

\begin{table}[!htbp]
    \centering
    \caption{Valoración de los factores del entorno.}
    \label{tab:ucp-fama-entorno}
    \small
    \setlength{\tabcolsep}{4pt}
    \renewcommand{\arraystretch}{1.10}
    \begin{tabular}{@{}lp{0.26\linewidth}p{0.40\linewidth}ccc@{}}
        \hline
        \textbf{ID} & \textbf{Factor} & \textbf{Justificación} &
        \textbf{Peso} & \textbf{Valor} & \textbf{Puntos} \\
        \hline
        F1 & Familiaridad con Objectory/RUP & La memoria documenta un proceso ágil basado en Scrum. & $1{,}5$ & 1 & $1{,}5$ \\
        F2 & Experiencia en aplicaciones similares & Valor neutro para la experiencia previa en sistemas MLOps acústicos. & $0{,}5$ & 3 & $1{,}5$ \\
        F3 & Experiencia en orientación a objetos & Valor neutro para la experiencia previa del equipo. & 1 & 3 & 3 \\
        F4 & Capacidad de análisis & Valor neutro para este factor del equipo. & $0{,}5$ & 3 & $1{,}5$ \\
        F5 & Motivación & Se infiere de la continuidad de las iteraciones y de los ajustes realizados tras las evaluaciones. & 1 & 4 & 4 \\
        F6 & Estabilidad de requisitos & Se conservaron los objetivos funcionales centrales durante las iteraciones. & 2 & 3 & 6 \\
        F7 & Dedicación parcial & Valor neutro ante la falta de un desglose de disponibilidad por integrante. & $-1$ & 3 & $-3$ \\
        F8 & Dificultad del lenguaje & Valor neutro para la dificultad de Python y TypeScript para el equipo. & $-1$ & 3 & $-3$ \\
        \hline
        \multicolumn{5}{@{}l}{\textbf{Total ponderado del entorno (EFactor)}} & $\mathbf{11{,}5}$ \\
        \hline
    \end{tabular}
\end{table}

\subsection{Cálculo de puntos y conversión a esfuerzo}

Los puntos sin ajustar (UUCP) se obtienen al sumar el peso de actores
y el de los casos de uso. Se aplican luego el factor de complejidad
técnica (TCF) y el factor del entorno (EF).

\[
    \mathrm{UUCP} = \mathrm{UAW} + \mathrm{UUCW} = 7 + 95 = 102
\]
\[
    \mathrm{TCF} = 0{,}6 + 0{,}01 \times 54 = 1{,}14
    \qquad
    \mathrm{EF} = 1{,}4 - 0{,}03 \times 11{,}5 = 1{,}055
\]
\[
    \mathrm{UCP} = \mathrm{UUCP} \times \mathrm{TCF} \times \mathrm{EF}
    = 102 \times 1{,}14 \times 1{,}055 \approx \mathbf{122{,}68}
\]

Según la regla de Schneider y Winters recogida en la pauta, se cuentan
los factores F1 a F6 con valoración inferior a 3 y los factores F7 a F8
con valoración superior a 3. Solo F1 cumple el criterio. Con un factor
desfavorable, la tasa de referencia es de 20 horas-persona por punto.
Al multiplicar el resultado sin redondear, se obtienen aproximadamente
\textbf{2\,454 horas-persona teóricas}.

La tasa de referencia se aplica como ejercicio metodológico. Su
productividad no está calibrada para el equipo ni para el alcance del
prototipo académico. El resumen de esfuerzo del apartado siguiente
presenta por separado las horas planificadas y registradas e incluye
la elaboración de la memoria.

% Fuentes del cálculo: capítulos 2, 5, 7 y 8 de la memoria actualizada.
% Factores y conversión a esfuerzo: pauta de estimación de la plantilla.
% Pesos de complejidad: Gustav Karner, Resource Estimation for Objectory
% Projects (1993), tablas 1 y 2.
